\begin{itemize}
        \item many attempts to understand how AI will affect work. but are we really measuring that correctly? No - existing benchmarks mainly measure automation. Most people aren't just using AI to do everything for you, most are using AI in a form of augmentation/assistance
        \item multi-agent systems, cost optimization
        \item some works measure assistance but on domain specific tasks - which LLMs are good for what domain? but not leaderboards on augmentation --> we want to understand augmentation and help others do that too - via a reproducible pipeline/framework
        \item are the top LLMs really the best for everything? often times the best players aren't the best coaches... the best researchers aren't the best teachers...
        \item a coach helping player vs the coach himself playing. a higher skill person helping a lower skill person
        \item our framework allows us to understand augmentation - helper vs helpee setting. 
        \item we not the first paper to compare auto vs augmentation but we are the first to evaluate them across seven professional tasks, with scalable measurement of role-specific capability profiles that standard agent benchmarks and human-only centaur studies do not isolate.
    \end{itemize}



% rough draft
With the rise in capabilities of artificial intelligence in recent years (or even months!), there have been numerous attempts to try to understand how this new-age Industrial Revolution will affect the daily lives of workers across the world. Organizations and individuals alike want to know how they will work with large language models (LLMs)-- if they will be able to incorporate them, or be replace them, or how costs may be affected. 

Studies across the domains of computer science, economics, and management have attempted to address these questions in the following ways. Computer scientists create large scale benchmarks of how LLMs can complete tasks across a number of domains and settings, evaluating and ranking their performances. But does this really get at how workers and organizations utilize LLMs in their daily work? This Stanford study (cite: https://futureofwork.saltlab.stanford.edu/) says no -- in fact the majority of workers prefer a collaborative partnership with AI. So if we are only measuring automation performance, we are missing a big part of the picture. A model which is the best at automating certain tasks is not necessarily the best at assisting, or augmenting, humans in completing those same tasks. One can refer to a famous sport analogy: the best players don't always make the best coaches.
    
Papers in the fields of economics and management have attempted to bridge this gap by testing the effects of human-AI collaboration in experimental settings. Yet thus far, each of these studies has focused on one specific task or domain (e.g. coding or working in a call center). Additionally, these papers usually test 1-2 specific models, which do not illustrate the full picture of how LLMs would perform in those settings.

What is missing in these bodies of work is a comparison of automation to augmentation. Fugener et al (CITE) illustrate this idea in a theoretical framework and verify their argument using an empirical study. [some other works also do this - add HERE] In this paper, we offer something new -- a scalable framework which evaluates LLMs across two important axes: (1) augmentation vs. automation and (2) multiple task domains. 

In the future, others can use this pipeline to understand how and which LLMs are best for their personal or organization's use case. By testing a task that is typical of their work, they an understand which LLM is the best fit for assisting them in their work, which is also going to be the most cost-efficient due to fit. Other organizations, who employ multi-agent systems, can also take this pipeline and apply it in their settings.